\documentclass[12pt]{article}
\usepackage[T1]{fontenc}
\usepackage[utf8]{inputenc}
\usepackage{lmodern}
\usepackage{textcomp}
\usepackage{enumitem}
\usepackage{geometry}
\usepackage{graphicx}
\usepackage{amsmath, amssymb}
\geometry{margin=1in}
\begin{document}
\begin{center}
History - K-12 Level Assessment 8 \
Total Marks: 40 \
Answer all questions.
\end{center}

\section*{Multiple Choice (10 marks)}
\begin{enumerate}[label=\arabic*.]
\item A celebration in which wealth is redistributed is called a:
\begin{enumerate}[label=(\alph*)]
    \item potlatch.
    \item kiva parties.
    \item totem feasts.
    \item sinagua.
\end{enumerate}
\item Subsistence practices of people can be determined by recovering and studying:
\begin{enumerate}[label=(\alph*)]
    \item the faunal assemblage from a site.
    \item an osteological comparative collection from a site.
    \item pumice and evidence of pyroclastic surges.
    \item all of the above.
\end{enumerate}
\item Which subfield of anthropology has been described as studying "other people's garbage"?
\begin{enumerate}[label=(\alph*)]
    \item primatology
    \item archaeology
    \item paleoanthropology
    \item linguistics
\end{enumerate}
\item Capacocha was an Incan ritual that involved what?
\begin{enumerate}[label=(\alph*)]
    \item drinking maize beer and feasting on the hearts of slain enemies
    \item sacrificing virgins in order to appease the serpent god
    \item sacrificing captives of war in order to propitiate the sun god
    \item sacrificing children in order to propitiate the spirits of the mountains
\end{enumerate}
\item When the ratio of brain size to body size is compared, which species has a proportionally larger brain?
\begin{enumerate}[label=(\alph*)]
    \item Homo erectus
    \item premodern Homo sapiens
    \item Neandertals
    \item modern Homo sapiens
\end{enumerate}
\end{enumerate}

\section*{True / False (10 marks)}
\textit{Answer True or False}
\begin{enumerate}[label=\arabic*.]
\setcounter{enumi}{5}
\item True or False: The city of Uruk and the cultures of the Maya and the Moche were all adversely affected by drought.
\item True or False: The undersea chasm between New Guinea/Australia and Java/Borneo is called the Bering Strait.
\item True or False: Homo sapiens and Australopithecus afarensis both rely on brachiation for locomotion, moving through trees like apes.
\item True or False: The three-age system describes human culture in a predictable chronology of the Stone Age, Bronze Age, and Iron Age.
\item True or False: By 1.8 million years ago, the hominid known as Homo erectus had evolved according to fossil evidence.
\end{enumerate}

\section*{Long Form Response (20 marks)}
\textit{Answer each question with a clear and complete explanation.}
\begin{enumerate}[label=\arabic*.]
\setcounter{enumi}{10}
\item Discuss the significance of the atlatl, an ancient spear-throwing device, in both Aztec culture and its broader historical context. How did it influence hunting and warfare?
\item Discuss the factors that contributed to New Zealand being one of the last regions occupied by human beings, and analyze the implications of this late settlement on its indigenous cultures.
\end{enumerate}

\vfill
\noindent\textit{End of Paper}
\end{document}